\subsubsection*{Объясните на графиках точки перегиба и асимптоты}

\textbf{Асимптоты.} На графиках коэффициента ускорения $K_y$ и эффективности $e$ асимптотами являются горизонтальные прямые, к которым кривая неограниченно приближается при возрастании аргумента ($r$ или $n$), но не достигает их на конечном участке. Значения асимптот получаются предельным переходом.

\textbf{Точки перегиба.} Точкой перегиба называется точка графика, в которой кривая меняет характер выпуклости (с «выгнутости вверх» на «выгнутость вниз» или наоборот). На графиках $K_y(r)$, $K_y(n)$, $e(r)$, $e(n)$ перегиб соответствует переходу от быстрого роста (или убывания) в начале к более пологому приближению к асимптоте.

\paragraph{График $K_y$ от ранга $r$ (при фиксированном $n$).}
Коэффициент ускорения задаётся формулой $K_y = \dfrac{r \cdot n}{n + r - 1}$. При $r \to +\infty$ выполняется
\[
\lim_{r \to +\infty} K_y = \lim_{r \to +\infty} \frac{r \cdot n}{n + r - 1} = n,
\]
поэтому \textbf{асимптотой} служит прямая $K_y = n$. С ростом $r$ кривая монотонно возрастает и прижимается к этой прямой снизу. \textbf{Точка перегиба} (если она одна) расположена на участке, где рост $K_y$ замедляется при приближении к асимптоте.

\paragraph{График $K_y$ от числа стадий $n$ (при фиксированном $r$).}
При $n \to +\infty$ имеем
\[
\lim_{n \to +\infty} K_y = \lim_{n \to +\infty} \frac{r \cdot n}{n + r - 1} = r,
\]
значит \textbf{асимптота} --- прямая $K_y = r$. Кривая $K_y(n)$ возрастает и приближается к ней. \textbf{Перегиб} соответствует зоне, где прирост ускорения с добавлением новых стадий становится всё меньше.

\paragraph{График $e$ от ранга $r$ (при фиксированном $n$).}
Эффективность $e = \dfrac{r}{n + r - 1}$. При $r \to +\infty$
\[
\lim_{r \to +\infty} e = \lim_{r \to +\infty} \frac{r}{n + r - 1} = 1,
\]
поэтому \textbf{асимптотой} является прямая $e = 1$. Кривая монотонно возрастает и стремится к единице снизу. \textbf{Точка перегиба} находится на участке перехода от крутого роста к пологому выходу на уровень $e = 1$.

\paragraph{График $e$ от числа стадий $n$ (при фиксированном $r$).}
При $n \to +\infty$
\[
\lim_{n \to +\infty} e = \lim_{n \to +\infty} \frac{r}{n + r - 1} = 0,
\]
значит \textbf{асимптота} --- ось абсцисс $e = 0$. Кривая $e(n)$ убывает и приближается к нулю. \textbf{Перегиб} соответствует области, где эффективность переходит от более крутого падения к более пологому стремлению к нулю.

Итого: на всех четырёх графиках горизонтальные асимптоты задаются указанными пределами; точки перегиба отвечают смене выпуклости при переходе от начального участка к асимптотическому поведению кривой.
